\documentclass[12pt]{article}

\usepackage{hyperref}
\usepackage{graphicx}
\usepackage{booktabs}
\usepackage{geometry}
\usepackage{tabularx}
\usepackage{float}
\usepackage[none]{hyphenat}
\usepackage{setspace}
\usepackage{subcaption}
\usepackage{titlesec}
\usepackage{amsmath}

\geometry{margin=1in}
\setstretch{0.96}

\titlespacing*{\section}{0pt}{10pt}{4pt}

\setlength{\parindent}{0pt}
\setlength{\parskip}{6pt}

\title{Labwork 1}
\author{Trinh Thanh Vinh - 2540110}
\date{\today}

\begin{document}

\maketitle

\section{Introduction}
In this labwork, we implement the gradient descent algorithm to find the minimum value of a simple quadratic function \( f(x) = x^2 \). We will analyze the convergence of the algorithm and observe how different learning rates affect the results.

\section{Implementation and Experimentation}
All the related functions and the gradient descent algorithm are implemented in the Python file \texttt{Labwork1.py} from scratch, without using any external libraries. The main components of the implementation include:
\begin{itemize}
    \item The function \( f(x) = x^2 \) which we want to minimize.
    \item An approximation of the derivative using the central difference method.
    \item The gradient descent algorithm that updates the value of \( x \) in each iteration, based on the computed gradient and a specified learning rate.
\end{itemize}

After building the implementation, we run the gradient descent algorithm with an initial value of \( x_0 = 10 \) and a learning rate of \( l = 0.1 \). The printed results include the step number, the value of \( x \) and \( f(x) \) in the current step, and the computed gradient. Different learning rates are also tested to observe their effects on convergence.

\section{Results and Discussion}

The results show that with a learning rate of \( l = 0.1 \), the minimum of \(x\) is reached at about $4.313 \times 10^{-7}$ after 76 steps. With smaller learning rates of 0.01 and 0.001, the algorithm takes more steps to converge, demonstrating the trade-off between learning rate and convergence speed. In particular, with \( l = 0.01 \), it takes 833 steps to reach a similar minimum value, while with \( l = 0.001 \), it takes 8398 steps.

On the other hand, when the learning rate is set at 0.5, the minimum value is reached in just 1 step, and the value of \( x \) is even closer to zero than with smaller learning rates. In fact, the algorithm still converges for learning rates between 0.6 and 0.9, but the number of steps it takes to reach the final value also increases as the learning rate grows, with 11 for \( l = 0.6 \), and 76 for \( l = 0.9 \). However, with a learning rate of 1, the value of the function eventually diverges to nearly 100, and the value of \(x\) jumps between 9 and -9, illustrating that too large a learning rate is not suitable for this problem.
\end{document}